\section{Related work}
\label{sec:related}

\subsection{Integrated scheduling of machines and vehicles}
\label{sec:related:plcsp}

Treating machine scheduling and vehicle dispatching as one problem has a
lineage that predates its current name. Ulusoy and
Bilge~\citep{ulusoy1993} scheduled machines and AGVs simultaneously and
proposed a genetic algorithm for the combined problem. The recent literature
has converged on the name production--logistics collaborative scheduling, used
by Li et al.~\citep{li2025plcsp} and Shi et al.~\citep{shi2025nested} among
others. The two names describe the same two-layer structure: a flexible job
shop decides which machine processes each operation, and a vehicle fleet
decides which vehicle performs each transfer.

The scope of the two layers is consistent across this literature, and so is one
boundary: reinforcement learning is used to control task allocation, that is
the assignment of operations to machines and of transfers to vehicles, while
local collision avoidance at the millisecond scale is delegated to the fleet
control system. We adopt the same boundary.

The problem also has an exact-optimisation line. Yao et al.~\citep{yao2026cp}
formulate the integrated problem as a constraint program and report that it
remains computationally intensive on the harder instances, which is the usual
motivation for learning-based approaches. The two lines are complementary
rather than competing.

Within the learning line, the decomposition varies. Hierarchical schemes
separate a high-level dispatching policy from lower-level execution, either as
a nested architecture~\citep{shi2025nested} or as an explicit layer
stack~\citep{wang2025hierarchical}; Wang et al.~\citep{wang2025hierarchical}
note that this design leans on the quality of the initial heuristic guidance
and can converge prematurely when that guidance is biased. Flat schemes place
all decisions in one policy, as Li et al.~\citep{li2025tsmc} do with a
multi-agent formulation over tasks, machines, and vehicles. Our formulation is
of the flat kind, and Section~\ref{sec:method:chain} gives the reason.

The set of resources being co-scheduled also varies. Yuan et
al.~\citep{yuan2026workers} schedule workers alongside AGVs in a distributed
flexible job shop, so that two heterogeneous transport resources compete for
the same jobs. We keep the resource set at machines and vehicles, which is the
boundary the PLCSP literature uses, and treat the wider resource set as out of
scope.

Whether the two layers are learned together or apart has itself become a
question. Link et al.~\citep{link2026coordination} compare joint and modular
training for a job shop with transport resources and find a coordination gap:
the modular arrangement does not recover what joint training obtains. That is
the closest published statement of the argument we make in
Section~\ref{sec:method:chain}, and it is about the training arrangement rather
than the architecture. Dong et al.~\citep{dong2026handling} reach the same
coupling from the exact-solution side, scheduling production together with
material-handling robots and carrying the robot battery state as an explicit
constraint. Dong et al.~\citep{dong2025heuristic} take the learning route on a
similar problem, with a heuristic supplying the dispatch decisions the policy
does not make.

\paragraph{Exact and metaheuristic treatments of the same coupling.}
A parallel line attacks the same production--logistics coupling with
hand-designed search rather than with a learned policy, and two of its members
are close enough to this paper to be worth stating precisely.
Chen et al.~\citep{chen2025carseq} integrate car sequencing on a mixed-model
assembly line with multi-trip vehicle routing and synchronized part supply; the
production sequence determines the due dates of the parts, which in turn
determine how the vehicles must be scheduled. That is the same asymmetric
causal chain we describe in Section~\ref{sec:method:chain}, in which production
creates the transport demand and fixes its deadlines. It is solved there
by a multi-start adaptive large-neighbourhood search with variable-neighbourhood
descent. Chen et al.~\citep{chen2026benders} couple order acceptance,
two-dimensional packing, and parallel machine scheduling, three decision layers
at once, and solve them by a nested logic-based Benders decomposition. Xu and
Xie~\citep{xu2025forest} couple scheduling with route planning in an emergency
setting, where the two layers also cannot be separated.

The methods in this line are problem-specific: an encoding, a decomposition, or
a set of cuts is designed per problem, and a change to the constraints
generally requires redesigning them. The same is true of the related machine
scheduling work in this line, where the algorithm is tailored to a
problem variant rather than transferred across
variants~\citep{chen2026upms}. The present work asks the complementary
question, whether a single learned policy over a state representation can serve
this class of coupled problems without per-problem algorithm design. We do not
claim that a learned policy beats these methods on the problems they were built
for; Section~\ref{sec:related:incomparable} explains why that comparison is not
available here, and Section~\ref{sec:results} reports where our policy does and
does not hold up against a metaheuristic on its own configuration.

\subsection{Learning architectures}
\label{sec:related:arch}

The dominant representation in recent work is a heterogeneous graph.
Xu et al.~\citep{taagnet2026} encode jobs, machines, and vehicles as three node
types with graph attention and type-aware pooling. Yang et
al.~\citep{yang2026hgamppo} combine heterogeneous graph attention with
multi-policy PPO. Yue et al.~\citep{yue2027hchgnn} build a hub-centric graph
for vehicle dispatching under finite-buffer constraints. Zhu et
al.~\citep{zhu2026dfjspht} apply a graph neural network to a distributed
flexible job shop with heterogeneous transport resources. Wu et
al.~\citep{wu2026madappo} use a heterogeneous graph attention network with
multi-head attention inside a two-agent PPO, for machine failures without
transport. Geng et al.~\citep{geng2025dtdrl} drive a multi-agent PPO from a
digital twin of the workshop, so the state is refreshed from a simulated copy
of the plant rather than from the shop itself.

A second line uses simpler sequence models. Li et al.~\citep{li2025plcsp} use a
multi-agent PPO with fully connected networks of 128-64-64 units and resolve
the final choice through a hand-written priority vector. Zhang et
al.~\citep{zhang2025sgformer} replace the graph encoder with a single-layer
simplified graph transformer to reach larger instances. Gao et
al.~\citep{gao2026gnntransformer} combine a graph convolutional stack with a
Transformer over task sequences. Wang et
al.~\citep{wang2025multitarget} train one policy to serve several scheduling
targets end to end, rather than one policy per target.

Our encoder is a standard Transformer~\citep{vaswani2017} over a token sequence
with no attention mask, and the novelty in this paper is not in the attention
mechanism. It is in what the attention is given to read, which is the token set
of Section~\ref{sec:method:state} and the transport-network-aware segment of
it. We note that the encoder choice is not a regression: recent work has moved
toward simplified representations and away from problem-specific graph
structure. Xiao et al.~\citep{xiao2026resched} reduce the flexible job shop to a
handful of state features behind a standard Transformer, and report that the
simplification does not cost performance; Wang et al.~\citep{wang2025memory} and
Zhang et al.~\citep{zhang2024improvement} keep a learned model in the loop but
use it to guide an improvement heuristic rather than to construct a schedule in
one pass. A plain attention encoder over a well-chosen token set sits within
that direction rather than against it.

Two further lines bear on how a policy is trained rather than what it reads.
Self-supervision can replace the reward signal entirely for job shop
scheduling~\citep{corsini2024selflabeling}, and a policy can be asked to serve
several preferences at once with self-evolution~\citep{choi2025multipolicy} or
to adapt per instance~\citep{qing2026instancewise}. Luttmann and
Xie~\citep{luttmann2026multiagentself} treat a set of simultaneous actions as
the object to improve, which is the closest formulation to our joint chain that
we have found outside this problem class. Stochastic processing times enter the
state through an attention encoder in~\citep{smit2025stochastic}. On the
application side, Zhu et al.~\citep{zhu2025matrix} couple the shop layout
itself to the scheduling policy, and Lee et
al.~\citep{lee2026digitaltwin} drive a dynamic AGV system from a digital twin.

What the representation must carry, however, is a substantive question and not
a matter of taste. Gao et al.~\citep{gao2026gnntransformer} state that their
encoder captures the static topology of the road network and therefore cannot
detect congestion arising from multi-vehicle interaction. Xu et
al.~\citep{taagnet2026} give vehicle features that contain neither coordinates
nor distances, so distance enters only through two hand-written dispatch rules.
Li et al.~\citep{li2025plcsp} use a network with no attention or graph
structure at all. In each case the constraint is absent from the state space,
and a policy cannot act on what its representation does not contain. The
mechanisms in Section~\ref{sec:method} that add aisle-segment tokens and a
distance bias to the attention score exist for this reason.

\subsection{Constraints and objectives}
\label{sec:related:constraints}

The constraints modelled in this literature vary, and the variation is
informative. Shi et al.~\citep{shi2025nested} state that machine breakdowns,
path conflicts, and AGV charging are ignored. Li et al.~\citep{li2025jms} state
that AGV failure is not considered and that time consumed by charging and
collision avoidance is negligible. Tang et
al.~\citep{tang2025lotstreaming} study static scheduling and do not consider
path conflict during travel. Wu et al.~\citep{wu2026madappo} state that energy
consumption is not involved and that workpieces and machines are all ready at
time zero, so transport does not enter. Zhang et
al.~\citep{zhang2026amr} treat maintenance of a deteriorating resource as
future work. Zhu et al.~\citep{zhu2026dfjspht} exclude transport route
conflicts and power replenishment. Yi et al.~\citep{yi2026petri} train without
accounting for AGV recharging needs. Li et al.~\citep{li2025plcsp} list AGV
charging, collision-avoidance-induced transport uncertainty, machine breakdown,
and worker absence as future work, and note that a DRL algorithm trained on one
workshop configuration needs retraining when the configuration changes.
Section~\ref{sec:related:coverage} revisits the same exclusions constraint by
constraint.



Charging is the exception to that pattern: papers do model battery state and
act on it, through a learned charging action~\citep{huang2025collaborative} or
a closed-form optimal rule~\citep{dong2026handling}, so a
claim that charging is unaddressed would be false. What is less often reported
is whether the constraint is ever active under the parameter values used:
whether a vehicle ever reaches the point where the decision matters.
Section~\ref{sec:results:active} measures that for our configuration, and finds
that it does not.

On objectives, the recent literature is predominantly makespan-first, with
energy entering either as a second objective or as a modelling contribution.
Yuan et al.~\citep{yuan2025green} build a five-component energy model with
multiple spindle speeds; Cheng et al.~\citep{cheng2025energy} combine deep
reinforcement learning with a memetic algorithm for an energy-saving variant.
Energy as a second objective also appears in the batch-processing line, where
Zheng et al.~\citep{zheng2025biobjective} solve a bi-objective non-identical
batch machine problem with makespan and energy exactly, and the batch
scheduling literature more generally is concerned with how batches are formed
and formulated~\citep{zheng2025formulations,xu2026autoclave,ding2024molding}. Energy is thus an established
objective, and the energy parameters we use are taken from Yuan et
al.~\citep{yuan2025green}. Multi-objective formulations appear in various
forms: a weighted sum over makespan and logistics
cost~\citep{shi2025nested}, and a Pareto-based hybrid that places a
non-dominated sorting step inside a DQN loop~\citep{tang2025lotstreaming}.

What we have not found in this line is a check on whether the reported
objectives are in fact independent of one another.
Section~\ref{sec:results:energy} runs that check. Preference conditioning is
the other current answer to the same problem: Smit et
al.~\citep{smit2026dcan} train a single policy over a preference vector and
recover a whole front from it, which is the natural extension of what this
paper reports for a fixed weight vector.

Maintenance is treated separately. Decisions about when to maintain or replace
a deteriorating system are studied in the reliability literature, where the
decision variable is a replacement or ordering policy rather than a scheduling
action~\citep{dong2024orderreplacement,dong2025sparepart}. Our maintenance
mechanism is of the scheduling kind (the policy chooses when to take an idle
machine out of service, within a forced threshold), and we are not aware of a
scheduling formulation of that decision in this problem class.

\subsection{How the ten constraints have been handled}
\label{sec:related:coverage}

Constraint coverage is set out below. The separation that matters is between a
constraint prior work does not handle and one it handles differently: not
having done something is weak evidence, whereas a treatment that differs in a
way that changes what the policy can do is a claim of substance.

\begin{description}\itemsep3pt
\item \textbf{C1 --- congestion.} \emph{Prior work.} Eleven works in this line exclude
      congestion or list it as future work. Three of them name the cause: the
      encoder carries only static topology, so ``the model cannot detect or
      avoid congestion that emerges in real time from multi-AGV
      interactions''~\citep{gao2026gnntransformer}; the vehicle features contain
      neither coordinates nor distances~\citep{taagnet2026}, and distance enters
      only through hand-written dispatch rules; and one policy network has no
      attention or graph structure at all~\citep{li2025plcsp}. Four works do
      choose a route, but leave congestion awareness to the simulator or to a
      rule. \emph{This paper.} The transport network enters the state: the aisle
      segments become tokens, a distance bias enters the attention score, and
      the route head reads the segments its candidate traverses
      (Section~\ref{sec:method:state}). \emph{Relation.} Not handled, with the
      cause stated by the authors. This is the strongest form of the claim
      available here, because it is a statement about the representation rather
      than about effort.
\item \textbf{C2 --- finite buffers.} \emph{Prior work.} Finite output buffers with a
      deadlock-aware priority mask that blocks all dispatching while any buffer
      holds a finished part~\citep{yue2027hchgnn}; others assume infinite
      capacity~\citep{taagnet2026,li2025plcsp}. \emph{This paper.} Buffer
      occupancy is a machine-token feature, and the constraint is otherwise left
      to the budget mechanism. \emph{Relation.} Handled differently: a
      hand-written mask versus a learned policy over the same state.
\item \textbf{C3 --- machine failure.} \emph{Prior work.} Memoryless failure at a fixed
      rate, and in several studies a single fixed failure scenario rather than
      a sampled process~\citep{wu2026madappo,geng2025dtdrl}. \emph{This paper.}
      The failure rate rises with the maintenance clock along a Weibull hazard.
      \emph{Relation.} Age dependence is not handled elsewhere. The rate law and
      its shape parameter are assumed, not calibrated to plant data.
\item \textbf{C4 --- rework.} \emph{Prior work.} A reworked operation returns to the task
      pool and is reassigned to a machine and a vehicle~\citep{li2025tsmc}.
      \emph{This paper.} Rework is performed in place on the same machine, as
      the simulation defines it, so no decision point is introduced.
      \emph{Relation.} Handled differently, by a modelling choice rather than by
      a mechanism.
\item \textbf{C5 --- setup.} \emph{Prior work.} Sequence-dependent setup as a
      machine--job edge attribute in the graph, with a shortest-setup-time
      baseline~\citep{yue2027hchgnn}; explicit setup tables in the general
      job-shop setting. Others assume immediate setup~\citep{taagnet2026}.
      \emph{This paper.} The setup time of a job change is a property of the
      pair rather than of either entity, so it is supplied to the machine head
      as a candidate feature. \emph{Relation.} Handled differently. Setup is the
      largest lever in our ablation (Table~\ref{tab:ablation}), so the treatment
      is consequential.
\item \textbf{C6 --- due dates.} \emph{Prior work.} Work-content due dates of the TWK
      family, with a per-job factor and an arrival term~\citep{taagnet2026,li2025plcsp};
      one work treats tardiness as a constraint rather than an
      objective~\citep{li2025plcsp}. \emph{This paper.} Due dates are exogenous,
      per job, and anchored on a load lower bound, so they do not move when the
      fleet configuration changes (Table~\ref{tab:duedates}). \emph{Relation.}
      Handled differently; exogeneity with respect to the fleet is the point of
      the design.
\item \textbf{C7 --- vehicle failure.} \emph{Prior work.} A failed vehicle releases its
      current task to the task pool and returns after
      repair~\citep{li2025plcsp}; three failure scenarios distinguished by when
      the failure occurs~\citep{shi2025nested}. \emph{This paper.} Tasks held by
      a failed vehicle are handed to vehicles that are still running, at
      checkpoints where the vehicle holds no aisle lock. \emph{Relation.}
      Handled differently: hand-over rather than return to a shared pool, with
      the target chosen by a rule rather than by the policy.
\item \textbf{C8 --- fleet and batching.} \emph{Prior work.} Vehicle capacity appears as
      a model parameter in five works, and none of them lets the policy decide
      how many tasks a trip should combine. Batching is a decision elsewhere,
      but a different one: in warehouse order batching the decision is which
      orders to pick together, coupled to the picker route rather than to a
      vehicle fleet~\citep{wang2025gphh}. We therefore do not claim that
      batching as a decision is new, only that transport batching inside a job
      shop is not where that decision has been studied. \emph{This paper.}
      Capacity is operational: one trip collects every queued task at a
      pick-up point and delivers to several drop points, and a separate head
      decides how many further tasks to append. \emph{Relation.} Handled
      differently, and the room the benchmark leaves for it is small
      (Table~\ref{tab:batching}).
\item \textbf{C9 --- charging.} \emph{Prior work.} Charging as a learned action over
      no-charge, opportunity-charge, and full-charge, with charger selection, on
      the same benchmark families~\citep{huang2025collaborative}; a closed-form
      optimal charging rule for a single-station variant~\citep{dong2026handling}.
      \emph{This paper.} A head chooses between not charging and each station
      when a vehicle is idle, and a vehicle at zero battery becomes unavailable
      until it recovers. \emph{Relation.} Handled differently, and not
      evaluated here: the constraint does not activate under the reported
      parameters (Table~\ref{tab:active}), so no claim is made about it.
\item \textbf{C10 --- maintenance.} \emph{Prior work.} No work in this line makes
      maintenance a decision. Four works list it as future work, and one models
      equipment deterioration while leaving the maintenance action
      open~\citep{zhang2026amr}. \emph{This paper.} A head chooses between
      starting maintenance on an idle machine now and deferring it, with the
      forced threshold as a hard floor. \emph{Relation.} Not handled elsewhere,
      but the constraint does not activate on the instance used for the policy
      results either, so the mechanism carries no measured effect here.
\end{description}

Two readings of this list are worth separating. The coverage claim itself: two
constraints are handled here as policy decisions that no surveyed work handles
that way, batch size and maintenance, and for congestion prior work not
only omits the constraint but states why its representation cannot carry it.
That last case is the substantive one, because it identifies a cause rather
than an omission.

The second reading is less comfortable. A constraint that no one has handled is
not automatically a contribution, and a constraint handled differently is a
difference rather than a novelty unless the difference changes what the policy
can do. For C2, C4, C5, C6, and C7 our treatment differs in where the
information enters, whether a candidate feature, a token, or a mask equivalent, and
these are differences of degree. For C8 and C10 the difference is categorical,
but the mechanisms carry no measured effect on the instance used here, so the
coverage is a statement about the model rather than about a demonstrated
result. What the list supports is that the model carries these constraints and
that the policy can act on them, not that each action improves the outcome.

\subsection{Where the field is moving}
\label{sec:related:landscape}

Four directions are visible in the most recent work on this problem, and our
design choices sit differently against each of them.

The heterogeneous graph with an attention encoder and a multi-policy PPO
backbone has become the default: it is the shared structure of the works by
Xu et al.~\citep{taagnet2026}, Yang et al.~\citep{yang2026hgamppo}, and Yue et
al.~\citep{yue2027hchgnn}. We depart from it in the encoder, for the
representation reasons given in Section~\ref{sec:related:arch}, and we keep the
single-policy structure for the advantage reasons given in
Section~\ref{sec:method:chain}.

Constraint realism is increasing. Finite buffers and deadlock, charging,
multi-trip load limits, deteriorating equipment, and worker availability now
appear individually in various papers~\citep{yue2027hchgnn,zhang2026amr,yuan2026workers}.
What is less common is a paper that carries many of them at once and reports
which of them are active. That is the gap this paper addresses, and
Section~\ref{sec:results:active} is the measurement.

Finally, the evaluation standard is rising. Recent work reports multiple
instance scales and, increasingly, real production data. Three current
benchmarks make the point: ATLAS draws on production traces from a large
manufacturing operator~\citep{jiang2026atlas}, DynaSchedBench calibrates
dynamic instances by difficulty and reports what it calls an observability
paradox in learned schedulers~\citep{cao2026dynaschedbench}, and FrontierCO
evaluates solvers at a scale that rules of thumb do not
reach~\citep{feng2026frontierco}. Our evaluation is the weakest part of this
paper by that standard, and Section~\ref{sec:discussion} takes it up.
